% !PRINTED 29
\weberpage{29}

The coefficients of the functions $Q$ and $C$ are composed from the
coefficients $a$ and $b$ by addition, subtraction, multiplication,
and division. In the denominator, however, only powers of $b_0$
occur; hence, if $b_0=1$, the coefficients of $Q$ and $C$ are whole
functions of the $a$ and $b$. The highest power of $b_0$ that can
occur in the denominator is the $(m-n+1)$st, since in the chain (4)
the factor $b_0$ is added once in the denominator in each succeeding
equation. In special cases, however, the highest power of $b_0$ in
all denominators may be lower.

Take, for example, for $A$ a function of the third degree,
\begin{equation}
 f(x)=a_0x^3+a_1x^2+a_2x+a_3
\tag{7}
\end{equation}
and for $B$ the so-called first derivative of $f(x)$, which is of the
second degree,
\begin{equation}
 f'(x)=3a_0x^2+2a_1x+a_2,
\tag{8}
\end{equation}
and hence:
\begin{equation}
 Q={1\over 3}x+{a_1\over 9a_0},
\tag{9}
\end{equation}
\begin{equation}
 C={6a_0a_2-2a_1^2\over 9a_0}x
   +{9a_0a_3-a_1a_2\over 9a_0}.
\tag{10}
\end{equation}

How the calculation prescribed in equations (4) is expediently
arranged may here be assumed as known from the elements. We draw
attention to the analogy between this calculation and division in the
decimal number system. A function $f(x)$ represents a number written
decimally if the coefficients $a_0,a_1,\ldots$ are whole numbers
between zero (inclusive) and $10$ (exclusive), and $x=10$ is set. If
numbers greater than $10$ are also permitted among the coefficients,
a number can be represented by $f(x)$ in different ways. This is used
in the division procedure to avoid fractional and negative
coefficients in the results in the simplest way.
